\documentclass[11pt, letterpaper]{article}

% Packages:
\usepackage[
    ignoreheadfoot, % set margins without considering header and footer
    top=2 cm, % seperation between body and page edge from the top
    bottom=2 cm, % seperation between body and page edge from the bottom
    left=1.5 cm, % seperation between body and page edge from the left
    right=1.5 cm, % separation between body and page edge from the right
    footskip=1.0 cm, % seperation between body and footer
    % showframe % for debugging 
]{geometry} % for adjusting page geometry
\usepackage{titlesec} % for customizing section titles
\usepackage{tabularx} % for making tables with fixed width columns
\usepackage{array} % tabularx requires this
\usepackage[dvipsnames]{xcolor} % for coloring text
\definecolor{primaryColor}{RGB}{0, 79, 144} % define primary color
\usepackage{enumitem} % for customizing lists
\usepackage{fontawesome5} % for using icons
\usepackage{amsmath} % for math
\usepackage[
    pdftitle={Mingyang Yao's CV},
    pdfauthor={Mingyang},
    pdfcreator={LaTeX with RenderCV},
    colorlinks=true,
    urlcolor=primaryColor
]{hyperref} % for links, metadata and bookmarks
\usepackage[pscoord]{eso-pic} % for floating text on the page
\usepackage{calc} % for calculating lengths
\usepackage{bookmark} % for bookmarks
\usepackage{lastpage} % for getting the total number of pages
\usepackage{changepage} % for one column entries (adjustwidth environment)
\usepackage{paracol} % for two and three column entries
\usepackage{ifthen} % for conditional statements
\usepackage{needspace} % for avoiding page brake right after the section title
\usepackage{iftex} % check if engine is pdflatex, xetex or luatex
\usepackage{etoolbox} % for \patchcmd
\usepackage{academicons} % Google Scholar icon (requires XeLaTeX or LuaLaTeX)
\usepackage{fontspec} % Required for academicons with XeLaTeX/LuaLaTeX


% Ensure that generate pdf is machine readable/ATS parsable:
\ifPDFTeX
    \input{glyphtounicode}
    \pdfgentounicode=1
    % \usepackage[T1]{fontenc} % this breaks sb2nov
    \usepackage[utf8]{inputenc}
    \usepackage{lmodern}
\fi



% Some settings:
\AtBeginEnvironment{adjustwidth}{\partopsep0pt} % remove space before adjustwidth environment
\pagestyle{empty} % no header or footer
\setcounter{secnumdepth}{0} % no section numbering
\setlength{\parindent}{0pt} % no indentation
\setlength{\topskip}{0pt} % no top skip
\setlength{\columnsep}{0cm} % set column seperation
\makeatletter
\let\ps@customFooterStyle\ps@plain % Copy the plain style to customFooterStyle
\patchcmd{\ps@customFooterStyle}{\thepage}{
    \color{gray}\textit{\small Mingyang Yao - Page \thepage{} of \pageref*{LastPage}}
}{}{} % replace number by desired string
\makeatother
\pagestyle{customFooterStyle}

\titleformat{\section}{\needspace{4\baselineskip}\bfseries\large}{}{0pt}{}[\vspace{1pt}\titlerule]

\titlespacing{\section}{
    % left space:
    -1pt
}{
    % top space:
    0.3 cm
}{
    % bottom space:
    0.2 cm
} % section title spacing

\renewcommand\labelitemi{$\circ$} % custom bullet points
\newenvironment{highlights}{
    \begin{itemize}[
        topsep=0.10 cm,
        parsep=0.10 cm,
        partopsep=0pt,
        itemsep=0pt,
        leftmargin=0.4 cm + 10pt
    ]
}{
    \end{itemize}
} % new environment for highlights

\newenvironment{highlightsforbulletentries}{
    \begin{itemize}[
        topsep=0.10 cm,
        parsep=0.10 cm,
        partopsep=0pt,
        itemsep=0pt,
        leftmargin=10pt
    ]
}{
    \end{itemize}
} % new environment for highlights for bullet entries


\newenvironment{onecolentry}{
    \begin{adjustwidth}{
        0.2 cm + 0.00001 cm
    }{
        0.2 cm + 0.00001 cm
    }
}{
    \end{adjustwidth}
} % new environment for one column entries

\newenvironment{twocolentry}[2][]{
    \onecolentry
    \def\secondColumn{#2}
    \setcolumnwidth{\fill, 4.5 cm}
    \begin{paracol}{2}
}{
    \switchcolumn \raggedleft \secondColumn
    \end{paracol}
    \endonecolentry
} % new environment for two column entries

\newenvironment{header}{
    \setlength{\topsep}{0pt}\par\kern\topsep\centering\linespread{1}
}{
    \par\kern\topsep
} % new environment for the header

\newcommand{\placelastupdatedtext}{% \placetextbox{<horizontal pos>}{<vertical pos>}{<stuff>}
  \AddToShipoutPictureFG*{% Add <stuff> to current page foreground
    \put(
        \LenToUnit{\paperwidth-2 cm-0.2 cm+0.05cm},
        \LenToUnit{\paperheight-1.0 cm}
    ){\vtop{{\null}\makebox[0pt][c]{
        \small\color{gray}\textit{Last updated in April 2026}\hspace{\widthof{Last updated in April 2026}}
    }}}%
  }%
}%

% save the original href command in a new command:
\let\hrefWithoutArrow\href

% new command for external links:
\renewcommand{\href}[2]{\hrefWithoutArrow{#1}{\ifthenelse{\equal{#2}{}}{ }{#2 }\raisebox{.15ex}{\footnotesize \faExternalLink*}}}


\begin{document}
    \newcommand{\AND}{\unskip
        \cleaders\copy\ANDbox\hskip\wd\ANDbox
        \ignorespaces
    }
    \newsavebox\ANDbox
    \sbox\ANDbox{}

    \placelastupdatedtext
    \begin{header}
        \textbf{\fontsize{20 pt}{20 pt}\selectfont Mingyang Yao}

        \vspace{0.1 cm}

        \normalsize
        \kern 0.25 cm%
        \mbox{\hrefWithoutArrow{mailto:m5yao@ucsd.edu}{\color{black}{\footnotesize\faEnvelope[regular]}\hspace*{0.13cm}m5yao@ucsd.edu}}%
        \kern 0.25 cm%
        \AND%
        \kern 0.25 cm%
        \mbox{\hrefWithoutArrow{tel:+1 858-305-1264}{\color{black}{\footnotesize\faPhone*}\hspace*{0.13cm}(858)305-1264}}%
        \kern 0.25 cm%
        \AND%
        \kern 0.25 cm%
        \mbox{\hrefWithoutArrow{https://scholar.google.com/citations?user=sBRCuJYAAAAJ&hl=en}{\color{black}{\aiGoogleScholar}\hspace*{0.15cm}Google Scholar}}%
        \kern 0.25 cm%
        \AND%
        \kern 0.25 cm%
        \mbox{\hrefWithoutArrow{https://andyweasley2004.github.io/}{\color{black}{\footnotesize\faGlobe}\hspace*{0.13cm}Personal Website}}%
    \end{header}

    \vspace{0.3 cm - 0.3 cm}

    \section{Education}
        \begin{twocolentry}{
            \textit{Aug 2026 --}}
            \textbf{University of Rochester}

            \textit{Ph.D in Electrical and Computer Engineering}

            \text{Supervisor: Prof.\ Zhiyao Duan}
        \end{twocolentry}

        \vspace{0.2 cm}

        \begin{twocolentry}{
            \textit{Sept 2022 -- Dec 2025}}
            \textbf{University of California, San Diego}

            \textit{BS in Mathematics-Computer Science}

            \textit{BS in Cognitive Science with spec.\ in Machine Learning and Neural Computation}

            \text{GPA: 3.93/4.00 -- Cum Laude}

            \text{Advisors: Dr.\ Ke Chen, Prof.\ Taylor Berg-Kirkpatrick, Prof.\ Shlomo Dubnov}
        \end{twocolentry}

    \section{Research Interests}
    \begin{onecolentry}
        \begin{highlightsforbulletentries}
            \item Music Generation

            \item Music Information Retrieval

            \item Music Representation Learning
        \end{highlightsforbulletentries}
    \end{onecolentry}


    \section{Publications}
    \begin{samepage}
        \begin{onecolentry}
            \textbf{From Generality to Mastery: Composer-Style Symbolic Music Generation via Large-Scale Pre-training \href{https://arxiv.org/abs/2506.17497}{Paper Link}}

            \vspace{0.10 cm}

            \mbox{\textbf{Mingyang Yao}}, \mbox{Ke Chen}

            \vspace{0.05 cm}

            \mbox{\textit{The Sixth Conference on AI Music Creativity (AIMC 2025); Oral Presentation}}
        \end{onecolentry}

        \vspace{0.2 cm}

        \begin{onecolentry}
            \textbf{BACHI: Boundary-Aware Symbolic Chord Recognition Through Masked Iterative Decoding on Pop and Classical Music \href{https://andyweasley2004.github.io/BACHI/}{Link}}

            \vspace{0.10 cm}

            \mbox{\textbf{Mingyang Yao}}, \mbox{Ke Chen}, \mbox{Shlomo Dubnov}, \mbox{Taylor Berg-Kirkpatrick}

            \vspace{0.05 cm}

            \mbox{\textit{IEEE International Conference on Acoustics, Speech, and Signal Processing (ICASSP 2026)}}

        \end{onecolentry}

        \vspace{0.2 cm}

        \begin{onecolentry}
            \textbf{The development of FEDUPP: Feeding Experimentation Device Users Processing Package to Assess Learning and Cognitive Flexibility \href{https://www.biorxiv.org/content/10.1101/2025.08.14.670424v1}{Paper Link}}

            \vspace{0.10 cm}

            \mbox{\textbf{Mingyang Yao$^*$}}, \mbox{Avraham M Libster$^*$}, \mbox{Shane Desfor}, \mbox{Freiya Malhotra}, \mbox{Nathalia Castorena},

            \mbox{Patricia Montilla-Perez}, \mbox{Francesca Telese}

            \vspace{0.05 cm}

            \mbox{\textit{Full Paper Submitted to Journal and Under Review}}

        \end{onecolentry}
    \end{samepage}


    \section{Research Experience}

        \begin{twocolentry}{
        \textit{Feb 2024 -- April 2025}}
            \textbf{Conditional Symbolic Music Generation}{ | Independent Research}

        \textit{Mentor: Ke Chen\href{https://www.knutchen.com/}{} | Adobe Research}
        \end{twocolentry}

        \vspace{0.10 cm}

        \begin{onecolentry}
        \begin{highlights}
            \item Developed 16k multi-genre MIDI {corpus}, 4 classical composer subsets and enhanced REMI representation
            \item Designed two-stage training combining large-scale pre-training with composer-specific fine-tuning, achieving superior few-shot style transfer compared to \href{https://huggingface.co/ElectricAlexis/NotaGen/blob/main/weights_notagen_pretrain-finetune_p_size_16_p_length_1024_p_layers_c_layers_6_20_h_size_1280_lr_1e-05_batch_1.pth}{NotaGen-finetuned}
            \item Established comprehensive evaluation framework integrating style fidelity metrics, harmonic analysis, and human preference studies to assess compositional quality
            \item Managed complete research pipeline, \href{https://generality-mastery.github.io/}{demo} page, and manuscript accepted at AIMC 2025
        \end{highlights}

        \end{onecolentry}

        \vspace{0.2 cm}

        \begin{twocolentry}{
        \textit{June 2025 -- Sept 2025}}
            \textbf{Symbolic Chord Recognition}{ | Independent Research}

        \textit{Mentor: Ke Chen\href{https://www.knutchen.com/}{}, Taylor Berg-Kirkpatrick | University of California San Diego}
        \end{twocolentry}

        \vspace{0.10 cm}

        \begin{onecolentry}
            \begin{highlights}
                \item Architected boundary-aware transformer with piano-roll patch embeddings, FiLM-based conditioning, and confidence-ordered masked decoding for chord prediction
                \item Demonstrated that confidence-ordered prediction achieves state-of-the-art performance with 6\% improvement on classical music benchmarks
                \item Curated POP909-CL dataset based on \href{https://github.com/music-x-lab/POP909-Dataset}{POP909} with expert annotations on chords and other metas, including time signature, key changes, and start beat
                \item Completed research pipeline, \href{https://andyweasley2004.github.io/BACHI/}{demo}, and manuscript accepted at ICASSP 2026
            \end{highlights}
        \end{onecolentry}

        \vspace{0.2 cm}

        \begin{twocolentry}{
        \textit{Sept 2023 -- Present}}
            \textbf{Telese Lab} {| Research Assistant}

        \textit{Principal Investigator: Francesca Telese}
        \end{twocolentry}

        \vspace{0.10 cm}
        \begin{onecolentry}
            \begin{highlights}
                \item Developed FEDUPP, the 1st open-source Python \href{https://github.com/ftlabucsd/FEDUPP}{pipeline} for automated FED3 feeding data analysis
                \item Integrated machine learning methods including unsupervised clustering for semi-automated labeling and LSTM classifiers for feeding quality assessment
                \item Led as co-first author, leading methodology design, experimental validation, and manuscript preparation with coordination among 4 people.
                \item Implemented computer vision pipeline using fine-tuned YOLOv11 and DeepLabCut for automated video processing and mouse trajectory tracking
            \end{highlights}
        \end{onecolentry}


    \section{Service}

    \textbf{Public Service}
    \vspace{0.1cm}

    \begin{twocolentry}{
    \textit{Feb 2026 -- Mar 2026}}
        \textbf{Program Chair}, IJCAI -- AI and Health Track 2026
    \end{twocolentry}

    \vspace{0.3 cm}

    \textbf{Teaching}
    \vspace{0.1cm}

    \begin{samepage}
    \begin{twocolentry}{
    \textit{Sept 2023 -- Dec 2023}}
        \textbf{Instructional Apprentice}

    \end{twocolentry}
        \begin{onecolentry}

        \textit{Course: Cogs 18 -- Introduction to Python | Instructor: Eric Morgan}
            \begin{highlights}
                \item Led weekly discussion sections and office hours, providing individualized guidance
                \item Graded assignments and exams with timely feedback, maintaining consistent standards across sections
                \item Achieved 100\% positive recommendation rate from student evaluations
            \end{highlights}
        \end{onecolentry}

        \vspace{0.2 cm}
        \begin{twocolentry}{
        \textit{April 2024 -- Dec 2024}}
            \textbf{Reader/Grader}

        \end{twocolentry}
        \begin{onecolentry}
            \vspace{0.1cm}
            \begin{highlights}
                \item Math 154 -- Discrete Math and Graph Theory (Spring 2024; Professor William Wesley)
                \item Math 20D -- Introduction to Differential Equations (Fall 2024; Professor Rishabh Dixit)
            \end{highlights}

        \end{onecolentry}
    \end{samepage}

    \vspace{0.1cm}
    \section{Awards and Honors}
    \begin{samepage}
        \begin{onecolentry}
            \text{UCSD Provost Honor (2022--2025)}

            \vspace{0.1cm}

            \text{Hefei No.1 High School International Department Salutatorian (2022)}

            \vspace{0.1cm}

            \text{Chinese Musicians' Association (CMA) Piano Graded Examination (Amateur), Level 10 (2015)}

        \end{onecolentry}
    \end{samepage}

\end{document}
